\chapter{2変数関数}
    \section{4象限逆正接 : \texorpdfstring{$\atanEx{x}{y}$}{atan2(y,x)}}
        \begin{shadebox}
            4象限逆正接 $\atanEx{x}{y}$ を次で定義する。
            \[
                \atanEx{x}{y} \coloneq
                \begin{cases}
                    \atan\parens*{\frac{y}{x}}+\pi & (\text{領域$1$ : }x<0,y>0)\\
                    \frac{\pi}{2} & (\text{領域$2$ : }x=0,y>0)\\
                    \atan\parens*{\frac{y}{x}} & (\text{領域$3$ : }x>0,y\neq0)\\
                    -\frac{\pi}{2} & (\text{領域$4$ : }x=0,y<0) \\
                    \atan\parens*{\frac{y}{x}}-\pi & (\text{領域$5$ : }x<0,y<0)
                \end{cases}
            \]
            $\atanEx{x}{y}$は原点及び$y$軸の負の領域の領域では定義されない。\\
        \end{shadebox}
        この関数は定義域で$C^1$級である。
        \begin{proof}
            \quad\par
            領域1,3,5では連続であり、偏微分可能であって偏微分係数はいずれも
            \begin{align*}
                \partDeriv{\atanEx{x}{y}}{x}{} &= \frac{-y}{x^2+y^2} \\
                \partDeriv{\atanEx{x}{y}}{y}{} &= \frac{x}{x^2+y^2}
            \end{align*}
            であり連続だから$C^1$級である。
            \par
            領域2における連続性を示す。領域2上の点$(a,b)\;(a=0,b>0)$に対して
            \[d = \left|\atanEx{a}{b}-\atanEx{a+h}{b+k}\ \right| = \left|\frac{\pi}{2}-\atanEx{b+k}{h}\right|\]
            とおくと、$k$を十分小さく取れば$b+k>0$なので
            \[
                d=
                \begin{cases}
                    \frac{\pi}{2}-\atan\parens*{\frac{b+k}{h}}\rightarrow0 & \text{as }h\rightarrow+0,k\rightarrow0 \\
                    \frac{\pi}{2}-\left(\atan\parens*{\frac{b+k}{h}}+\pi\right) & \text{as }h\rightarrow-0,k\rightarrow0 \\
                \end{cases}
            \]
            であり、領域2における連続性が示せた。領域4における連続性も同様に示せる。
            \par
            さて、先程の偏微分係数をみると、幸運なことにこれは領域2,4においても連続である。
            以上より結局、領域1,2,3,4,5全域で$\atanEx{x}{y}$は$C^1$級である。
        \end{proof}